\documentclass[dvisvgm,tikz,border=5pt]{standalone}
\usepackage{amsmath,amssymb}
\usepackage{CJKutf8}
\usetikzlibrary{arrows.meta,positioning,calc}

\definecolor{themebg}{HTML}{2e362c}
\definecolor{themefg}{HTML}{edf4e8}
\definecolor{themegray}{HTML}{9ea897}
\definecolor{themecurve}{HTML}{7eb6ff}
\definecolor{themealert}{HTML}{ff7b7b}
\definecolor{themeamber}{HTML}{f5b942}

\pgfdeclarelayer{background}
\pgfdeclarelayer{foreground}
\pgfsetlayers{background,main,foreground}

\begin{document}
\begin{CJK*}{UTF8}{gbsn}
\pagecolor{themebg}
\begin{tikzpicture}[>=Stealth,font=\small,color=themefg,text=themefg]
\tikzset{
  bag/.style={rounded corners=4pt,draw=themegray,fill=themefg!3!themebg,minimum height=.9cm,inner xsep=7pt,inner ysep=5pt,align=center},
  act/.style={rounded corners=4pt,draw=themeamber,fill=themeamber!10!themebg,minimum height=.9cm,inner xsep=8pt,align=center},
  edge/.style={->,themegray,line width=1pt}
}

\node[anchor=west,font=\bfseries] at (0,5.9)
  {4 叉 Huffman：补零与“首轮少合并”只是两种等价记账方式};
\node[anchor=west,font=\small,text=themegray] at (0,5.35)
  {真实权值 $\{1,1,2,3,5,8,13,21\}$，因为 $8\not\equiv1\pmod3$，需要补两个零权虚叶。};

% left route
\node[font=\bfseries,text=themecurve] at (4.0,4.45) {视角 A：显式补两个 0};
\node[bag] (a0) at (4.0,3.55) {$\{0,0,1,1,2,3,5,8,13,21\}$};
\node[act] (a1) at (4.0,2.25) {$0+0+1+1=2$};
\node[bag,draw=themecurve] (a2) at (4.0,.95) {$\{2,2,3,5,8,13,21\}$};
\draw[edge] (a0)--(a1); \draw[edge] (a1)--(a2);

% right route
\node[font=\bfseries,text=themeamber] at (11.2,4.45) {视角 B：不画虚叶};
\node[bag] (b0) at (11.2,3.55) {$\{1,1,2,3,5,8,13,21\}$};
\node[act] (b1) at (11.2,2.25) {首轮只取两个：$1+1=2$};
\node[bag,draw=themecurve] (b2) at (11.2,.95) {$\{2,2,3,5,8,13,21\}$};
\draw[edge] (b0)--(b1); \draw[edge] (b1)--(b2);

\draw[<->,themecurve,line width=1.1pt] (a2.east)--node[above,font=\small,text=themecurve]{同一候选多重集}(b2.west);
\node[anchor=west,font=\small,text=themegray] at (1.8,-.25)
  {从这里开始，两种写法完全一致：之后每一轮都取当前 4 个最小权值合并。};

\end{tikzpicture}
\end{CJK*}
\end{document}
